\documentclass[11pt]{article}
\usepackage[margin=1in]{geometry}
\usepackage{amsmath}
\usepackage{amssymb}
\usepackage{hyperref}
\usepackage[numbers]{natbib}
\usepackage{booktabs}
\usepackage{multirow}

\hypersetup{
    colorlinks=true,
    linkcolor=blue,
    citecolor=blue,
    urlcolor=blue
}

\title{Daily Paper Review --- March 21, 2026\\
\large CubiD: Cubic Discrete Diffusion on\\High-Dimensional Representation Tokens}
\author{Internbot --- Automated Research Review}
\date{March 21, 2026}

\begin{document}
\maketitle

\begin{abstract}
This report reviews \textbf{CubiD}~\cite{cubid}, the first discrete diffusion model that generates directly on high-dimensional (768-dim) representation tokens from pretrained vision encoders (DINOv2, SigLIP2).
Current discrete generation methods are confined to low-dimensional tokens (8--32 dims), which sacrifice the semantic richness needed for understanding tasks.
CubiD's key innovation is \textbf{per-element masking}: treating the $h \times w \times d$ representation tensor as a 3D ``cube'' where any individual scalar can be independently masked, forcing the model to learn both within-position dimensional correlations and across-position spatial correlations.
This achieves state-of-the-art gFID of 1.88 on ImageNet-256 among all discrete methods, while preserving $>$99\% of the encoder's understanding capability.
We connect this to Omni-Diffusion~\cite{omnidiff} (masked discrete diffusion for multimodal generation): both leverage masked discrete diffusion, but CubiD extends it to high-dimensional token spaces, opening a path toward unified multimodal architectures where text and vision share semantically meaningful discrete tokens.
\end{abstract}

\section{Paper Summary}

\subsection{Overview}
\begin{itemize}
    \item \textbf{Title:} Cubic Discrete Diffusion: Discrete Visual Generation on High-Dimensional Representation Tokens
    \item \textbf{Authors:} Yuqing Wang, Chuofan Ma, Zhijie Lin, Yao Teng, Lijun Yu, Shuai Wang, Jiaming Han, Jiashi Feng, Yi Jiang, Xihui Liu (HKU, ByteDance Seed, CMU, NJU, CUHK)
    \item \textbf{arXiv:} \href{https://arxiv.org/abs/2603.19232}{2603.19232} \quad (CVPR 2026)
    \item \textbf{Code:} \href{https://github.com/YuqingWang1029/CubiD}{github.com/YuqingWang1029/CubiD}
    \item \textbf{Topics:} Discrete diffusion, high-dimensional token generation, unified representation
\end{itemize}

\subsection{Core Problem}

Pretrained vision encoders (DINOv2, SigLIP2) produce high-dimensional representation tokens (768--1024 dims) that excel at understanding tasks.
But current discrete generation methods cannot operate on these tokens:
\begin{itemize}
    \item \textbf{Vector quantization fails}: A 768-dim token quantized to $L$ levels per dimension yields a codebook of size $L^{768}$---combinatorially intractable.
    \item \textbf{Standard masked diffusion} (MaskGIT) treats each spatial position as an atomic unit from a finite vocabulary. This works for 8--32 dim tokens but \emph{cannot learn within-position correlations} at 768 dims.
    \item \textbf{Autoregressive generation} would require $h \times w \times d = 16 \times 16 \times 768 = 196{,}608$ sequential steps---prohibitively slow.
\end{itemize}

The result: generation and understanding use different token spaces, preventing unified multimodal architectures.

\subsection{Proposed Method}

\paragraph{Dimension-Wise Scalar Quantization.}
Instead of vector quantization, CubiD independently quantizes each of the 768 feature dimensions into $L$ discrete levels:
\[
q_{x,y,i} = \mathrm{Quantize}(z_{x,y,i};\, L)
\]
This decomposes the intractable $L^{768}$ codebook into 768 independent $L$-level quantizers.
Optimal levels: $L\!=\!8$ for DINOv2-B (rFID 0.57), $L\!=\!16$ for SigLIP2-B (rFID 0.69).
Critically, dimension-wise quantization preserves $>$99\% of understanding capability (GQA: 63.1 vs.\ 63.2 continuous; TextVQA: 59.8 vs.\ 59.6), while vector quantization degrades catastrophically (GQA drops to 54.9).

\paragraph{Per-Element (Fine-Grained) Masking.}
The representation tensor has shape $h \times w \times d$ (e.g., $16 \times 16 \times 768$).
CubiD treats this as a 3D ``cube'' and masks at the \textbf{individual scalar level}---any element at any spatial position and any feature dimension can be independently masked.

\begin{center}
\begin{tabular}{lrl}
\toprule
\textbf{Masking Strategy} & \textbf{gFID} & \textbf{Failure Mode} \\
\midrule
Per-dimension (entire dim across all positions) & 120.03 & Texture artifacts, no spatial coherence \\
Per-spatial (entire vector at a position) & 22.22 & Blurry, locally inconsistent \\
\textbf{Per-element (CubiD)} & \textbf{5.33} & High quality \\
\bottomrule
\end{tabular}
\end{center}

Per-element masking is \emph{categorically necessary}: it forces the model to learn (1) within-position dimensional correlations (given dims 1--400 at position $(3,5)$, predict dims 401--768), (2) across-position spatial correlations (standard spatial coherence), and (3) cross-axis correlations spanning both axes simultaneously.

\paragraph{Architecture.}
A bidirectional Transformer processes $h \times w = 256$ spatial tokens, each of dimension $d\!=\!768$.
The prediction head produces $d \times L$ logits per position, predicting all 768 dimensions simultaneously with $L$-way classification per dimension.
Model sizes: CubiD-L (946M), CubiD-XL (1.4B), CubiD-XXL (3.7B).

\paragraph{Training.}
Mask ratio $r$ is sampled from a truncated Gaussian biased toward aggressive masking: $r \sim \mathrm{TruncNorm}(\mu\!=\!1.0, \sigma\!=\!0.10, [0, 1])$.
Loss is standard cross-entropy over masked positions:
\[
\mathcal{L} = -\mathbb{E}_{q, M}\left[\sum_{i \in M} \log p(q_i \mid q_{\bar{M}})\right]
\]
A learnable [MASK] embedding replaces masked positions (using random codebook samples instead yields gFID 56.38 vs.\ 5.33).

\paragraph{Fixed-Step Generation.}
Inference proceeds iteratively from a fully masked tensor: at each step, the model predicts all masked elements simultaneously, a subset is accepted (cosine schedule), and the rest remain masked.
The number of steps $T$ is \textbf{independent of feature dimensionality $d$}---whether tokens are 32-dim or 768-dim, the model requires the same $\sim$256--512 steps.
This is because the Transformer processes $h \times w$ spatial tokens (not $h \times w \times d$ individual elements), and each forward pass predicts all $d$ dimensions at every position.

\subsection{Key Results}

\paragraph{State-of-the-Art Among Discrete Methods.}
On ImageNet-256 (class-conditional):

\begin{center}
\begin{tabular}{llccc}
\toprule
\textbf{Method} & \textbf{Type} & \textbf{Token Dim} & \textbf{Params} & \textbf{gFID} \\
\midrule
MaskGIT & Masked diffusion & 16 & 227M & 4.02 \\
VQ-Diffusion & Discrete diffusion & 16 & 370M & 11.89 \\
LlamaGen-XXL & Autoregressive & 8 & 1.4B & 2.34 \\
VAR & Autoregressive & 32 & 2.0B & 1.97 \\
VFMTok-3B & AR (reorganized repr.) & 12 & 3.1B & 2.04 \\
\textbf{CubiD-XXL} & \textbf{Masked diffusion} & \textbf{768} & \textbf{3.7B} & \textbf{1.88} \\
\bottomrule
\end{tabular}
\end{center}

CubiD achieves the best gFID among all discrete generation methods, operating at 768 dimensions---24$\times$ higher than the next best.

\paragraph{Understanding Preservation.}
When used in LLaVA for visual understanding:

\begin{center}
\begin{tabular}{lcccc}
\toprule
\textbf{Tokenizer} & \textbf{GQA} & \textbf{TextVQA} & \textbf{POPE} & \textbf{MME} \\
\midrule
SigLIP2 (continuous) & 63.2 & 59.6 & 85.4 & 1484 \\
SigLIP2-DQ (CubiD) & 63.1 & 59.8 & 85.0 & 1480 \\
SigLIP2-VQ (vector quant.) & 54.9 & 45.6 & 81.2 & 1189 \\
\bottomrule
\end{tabular}
\end{center}

Dimension-wise quantization preserves understanding nearly perfectly; vector quantization destroys it.

\paragraph{Scaling.}
Consistent improvement from 946M to 3.7B parameters (gFID: 5.25 $\rightarrow$ 4.91 $\rightarrow$ 4.68 without cfg).
The model also generalizes to low-dimensional tokens: on ImageNet-512 with 32-dim DC-AE tokens, CubiD achieves gFID 1.58 (0.95B params), outperforming USiT-2B (1.72, 1.58B params).

\paragraph{Ablations.}
Additional design choices validated: (1) mask ratio $\sigma\!=\!0.10$ is optimal; (2) DINOv2-B slightly outperforms SigLIP2-B for generation (gFID 5.25 vs.\ 5.87); (3) generation steps saturate at $\sim$512; (4) compressing 768-dim to 32-dim before generation improves gFID (1.55) but degrades understanding---CubiD deliberately preserves the full dimensionality for dual-use capability.

\subsection{Limitations}
\begin{itemize}
    \item \textbf{Decoder dependence}: Reconstruction quality is capped at $\sim$18dB PSNR by the RAE decoder, limiting fine-grained visual details.
    \item \textbf{Generation--reconstruction gap}: A gap remains between CubiD and continuous diffusion methods using the same autoencoders.
    \item \textbf{Inference speed}: Requires 256--512 steps per image. Faster than autoregressive (196K steps) but slower than some continuous diffusion methods.
    \item \textbf{Vision-only evaluation}: The unified generation--understanding capability is validated but not demonstrated in a full multimodal system.
\end{itemize}

\section{Follow-up Research Directions}

\subsection{Per-Element Masking for Diffusion Language Models}

\paragraph{Gap.}
Omni-Diffusion~\cite{omnidiff} uses masked discrete diffusion for text generation, but treats each text token as an atomic unit from a fixed vocabulary (per-spatial masking in CubiD's terminology).
CubiD shows that per-element masking is categorically better (gFID 5.33 vs.\ 22.22) because it captures within-position correlations.
For language, this is unexplored: standard text tokens are 1D codebook indices, but recent work on multi-attribute tokenization (separate semantic, syntactic, and phonetic codes per position) creates a natural ``cube'' structure where per-element masking could apply.

\paragraph{Approach.}
Design a language tokenizer that produces multi-dimensional discrete codes per position (e.g., a $T \times d$ tensor where $T$ is sequence length and $d$ encodes separate linguistic attributes).
Apply CubiD's per-element masking during training: independently mask individual attribute dimensions at each position.
This would force the model to learn correlations between linguistic attributes (e.g., given the semantic code, predict the syntactic code) and across positions (standard language modeling).
Omni-Diffusion's multimodal framework could serve as the backbone, replacing its per-token masking with per-element masking.

\paragraph{Feasibility.}
The architecture transfers directly: a bidirectional Transformer over $T$ positions with $d \times L$ logits per position.
The main challenge is designing the multi-dimensional text tokenizer---residual VQ (RVQ) provides a natural starting point, where each VQ level becomes a ``dimension.''
Initial experiments could use Omni-Diffusion's text generation benchmarks to compare per-token vs.\ per-element masking on RVQ-encoded text.

\subsection{Efficient CubiD via IndexCache-Style Cross-Layer Reuse}

\paragraph{Gap.}
CubiD requires 256--512 iterative demasking steps, each involving a full Transformer forward pass.
IndexCache~\cite{indexcache} showed that adjacent Transformer layers share 70--100\% of their sparse attention indices, enabling cross-layer index reuse with 1.82$\times$ speedup.
CubiD's bidirectional Transformer has the same layer-wise redundancy, but no work has applied cross-layer computation reuse to iterative discrete diffusion.

\paragraph{Approach.}
Apply IndexCache's F/S (Full/Shared) layer partitioning to CubiD's generation Transformer: Full layers compute fresh attention; Shared layers reuse cached indices from the previous Full layer.
Additionally, exploit \textbf{cross-step} reuse: between consecutive demasking steps, only a fraction of tokens change (the newly unmasked positions).
Cache attention computations for unchanged positions and only recompute for modified ones, similar to KV cache reuse in autoregressive models.
This double reuse (cross-layer + cross-step) could provide multiplicative speedups.

\paragraph{Feasibility.}
IndexCache is training-free in its basic variant and applies to any Transformer backbone.
CubiD's cosine demasking schedule means early steps unmask many tokens (high change rate) while late steps unmask few (low change rate)---cross-step reuse becomes increasingly effective as generation progresses.
A 2$\times$ speedup from cross-layer reuse combined with 1.5--2$\times$ from cross-step caching could reduce 512-step generation to an effective $\sim$128--170 steps of full computation.

\subsection{CubiD for Unified Vision-Language Discrete Diffusion}

\paragraph{Gap.}
CubiD operates on vision tokens and Omni-Diffusion operates on text tokens, but neither handles both modalities in a shared high-dimensional discrete space.
CubiD's dimension-wise quantization preserves understanding ($>$99\%), and Omni-Diffusion's masked diffusion handles text generation.
A unified model that applies CubiD's per-element masking to \emph{both} vision and language tokens in a shared representation space would enable true multimodal discrete diffusion.

\paragraph{Approach.}
Train a shared encoder that maps both text and images into a common $d$-dimensional discrete representation space (e.g., $d\!=\!768$).
For vision: use CubiD's dimension-wise quantization of DINOv2/SigLIP2 features.
For language: train a text encoder that produces 768-dim discrete tokens aligned with the vision space (via contrastive or distillation objectives).
Apply CubiD's per-element masked diffusion over the concatenated vision-language token sequence, enabling the model to jointly generate and understand both modalities.
The per-element masking would capture cross-modal correlations: given partial vision features and partial text features, predict the missing elements of both.

\paragraph{Feasibility.}
CubiD's architecture is modality-agnostic---the bidirectional Transformer processes sequences of $d$-dimensional tokens regardless of whether they originate from vision or language.
The main challenge is training the shared discrete representation space to preserve both generation quality and understanding capability across modalities.
SigLIP2 already aligns vision and language in continuous space; extending this alignment to the discrete quantized space is a natural next step.
Evaluation could use Omni-Diffusion's multimodal benchmarks (image-text generation, VQA, captioning).

\section{Conclusion}

CubiD demonstrates that discrete diffusion can scale to high-dimensional (768-dim) representation tokens through two key innovations: dimension-wise scalar quantization (decomposing the intractable codebook) and per-element fine-grained masking (capturing both within-position and across-position correlations).
The result is the first discrete generation model that achieves state-of-the-art generation quality (gFID 1.88) while preserving $>$99\% of the encoder's understanding capability---a critical step toward unified multimodal architectures.

This paper pairs naturally with Omni-Diffusion in our review series: both leverage masked discrete diffusion, but at different granularities.
Omni-Diffusion masks entire tokens (per-spatial); CubiD masks individual scalar elements (per-element), achieving a 4$\times$ quality improvement.
The synthesis is clear: applying CubiD's per-element masking to Omni-Diffusion's multimodal framework could enable higher-quality text and image generation in a shared high-dimensional discrete space.
Combined with IndexCache for inference acceleration, this outlines a path toward efficient, unified discrete diffusion for both language and vision.

\begin{thebibliography}{9}
\bibitem{cubid} Wang, Y., Ma, C., Lin, Z., Teng, Y., Yu, L., Wang, S., Han, J., Feng, J., Jiang, Y., \& Liu, X. (2026). Cubic Discrete Diffusion: Discrete Visual Generation on High-Dimensional Representation Tokens. \textit{CVPR 2026}. arXiv:2603.19232.

\bibitem{omnidiff} (2026). Omni-Diffusion: Unified Multimodal Understanding and Generation with Masked Discrete Diffusion. \textit{arXiv preprint arXiv:2603.06577}.

\bibitem{indexcache} Bai, Y., et al. (2026). IndexCache: Accelerating Sparse Attention via Cross-Layer Index Reuse. \textit{arXiv preprint arXiv:2603.12201}.
\end{thebibliography}

\end{document}
